\section{Technical Overview}

Here, we provide the background and exposition of the main ideas in our work.

\subsection{The upper bound}

\paragraph{Classical Krivine rounding schemes.}
A classical Krivine scheme consists of odd sign functions \(f,g:\mathbb R^k\to\{-1,1\}\). Its normalized correlation function is
\(
    H_{f,g}(t):=\frac{\pi}{2}\mathbb E[f(X)g(Y)],
\)
where \(X,Y\) are standard Gaussian vectors with coordinatewise correlation \(t\). Suppose \(H^{-1}(z)=\sum_{n\geq 1}a_nz^n\), and let its \emph{majorant} series be  \(M(s):=\sum_{n\geq 1}|a_n|s^n\). Krivine's tensor preprocessing
(\cref{app:krivine-preprocessing})
gives
\(
    M(\gamma)\leq 1
    \implies 
    K_G\leq\frac{\pi}{2\gamma}.
\)
The upper-bound problem is therefore to construct a correlation function whose inverse has a large admissible  $\gamma$ so that $M(\gamma) = 1$.


\paragraph{Limiting Krivine rounding schemes.}
In ordinary Krivine schemes, we need individual Gaussian coordinates $X_i, Y_i$ to have coordinate-wise correlation $t$. We enlarge the search space by allowing $X_i, Y_i$ to have coordinate-wise correlation
\[
    \rho_i(t)=\sum_{\substack{d\geq 1\\ d\ \mathrm{odd}}}c_{i,d}t^d,
    \qquad
    \sum_d|c_{i,d}|\leq 1.
\]
This gives us more degrees of freedom in designing \(H\) whose associated \(M\) has larger \(\gamma\) satisfying $M(\gamma) = 1$.

% This flexibility in correlation structure allows us to find correlation functions \(H\) whose corresponding inverse majorants \(M\) have larger radii \(\gamma\).

% In our construction, we use the central limit theorem to realize arbitrary correlations as above.  
% In our construction, the powers \(t^d\) arise from Hermite polynomials: if \(U,V\) are correlated standard Gaussians, then \(\mathbb E[\operatorname{He}_d(U)\operatorname{He}_d(V)]=t^d\). 

% A limiting scheme is thus a fixed-dimensional model of a family of high-dimensional schemes built from normalized sums of Hermite polynomials.

% \rnote{The previous sentence is too vague and terse. You can just state the idea; we come up with an intermediate function and use central-limit theorem to realize arbitrary correlations as above.}


\paragraph{Approximation of limiting schemes using classical schemes.}
% While we may design the correlation function \(H\) to admit any \(\rho\) of the above structure, it is not immediately 
While we may design the correlations between variables to admit any \(\rho\) of the above structure, it is not immediately clear how the rounding algorithm works. 
Here, we show that any limiting scheme can in fact be approximated arbitrarily well by classical Krivine schemes, whose rounding algorithm is established. 
Therefore, it suffices to find a limiting scheme whose inverse majorant $M$ admits a high $\gamma$ with $M(\gamma) = 1$.
% In order to use limiting schemes, we show that every limiting scheme can be approximated by classical Krivine schemes.

More precisely, we construct a sequence $\mathcal{S}_1, \mathcal{S}_2, \ldots$ of classical Krivine schemes of increasing dimension \(k\), whose correlation functions $H_{\mathcal{S}_1}, H_{\mathcal{S}_2}, \dots$ converge to the correlation function $H_{\mathcal{S}}$ of a limiting scheme $\mathcal{S}$.
The main idea is to use the central limit theorem to collapse the increasing number of Gaussian random variables in the classical schemes \(\mathcal{S}_k = (f_k, g_k)\), where \(f_k, g_k : \mathbb{R}^k \to \{\pm 1\}\).

  
% Therefore, it suffices to find a limiting rounding scheme that gives a good upper bound on \(K_G\), and a way to generate a sequence of classical rounding schemes that asymptotically approximates it.

% The main idea of the proof is to use the central limit theorem to collapse the increasing dimensions.

% This lets us search for the upper-bound construction entirely within the class of limiting schemes.

% \rnote{Reads a bit like bullet points. No need for putting Vitali's theorem here. Need less jargon here.}


% 

\paragraph{The cubic--quintic scheme.}
A limiting Krivine scheme that attains the paper's stated upper bound on \(K_G\) is defined as follows.
% Our construction is a two-dimensional limiting scheme.
For suitable parameters \(s_3,s_5,\vartheta>0\), set
\[
    \rho(t):=
    \frac{t-s_3^2t^3+s_5^2t^5}{1+s_3^2+s_5^2},
\]
and use the partitions
\[
    f(w,x)=\operatorname{sgn}\!\left(w+\vartheta\operatorname{He}_3(x)\right),
    \qquad
    g(w,x)=\operatorname{sgn}\!\left(w-\vartheta\operatorname{He}_3(x)\right).
\]
We set the first coordinate \(w\) to have correlation \(\rho(t)\), and set the second coordinate \(x\) to have correlation \(t\).

% This limiting Krivine scheme gives our stated upper bound on \(K_G\).
% , and we show how to approximate it by classical schemes.

% The first Gaussian coordinate pair has correlation \(\rho(t)\), while the second has correlation \(t\). 
% The cubic direction is anti-aligned between the two sides and the quintic direction is aligned. 
% The parameters are chosen to cancel out the first nonlinear coefficients of the resulting correlation function. 
% \rnote{Could explicitly spell out, the final form $H_{f,g}$.}


\paragraph{Deriving the upper bound.}
In order to prove the upper bound (which we denote as $K_G^{UB}$), one needs to show that the inverse majorant $M(\gamma) \leq 1$ holds, for $\gamma=\frac{\pi}{2K_G^{UB}}$.
Doing this is non-trivial, and requires both an analytic step and a computational step. 

For the analytic step, write
\(H(t)=b_1t+\sum_{m\geq 3}b_mt^m\), and let
\(\Delta_H:=\sum_{m\geq 3}|b_m|\) denote its total nonlinear coefficient mass. We prove that it is sufficient to show 
\[
    \gamma+\Delta_H<b_1.
\]
Thus, rather than computing the inverse series \(H^{-1}\), it is enough to certify a lower bound on the linear coefficient \(b_1\) and an upper bound on \(\Delta_H\).
We then compute the first terms of $\{b_n\}$ using a formula involving the Hermite coefficients of $f$ and $g$, which are one-dimensional Gaussian integrals. Finally, we bound the tail of $\Delta_H$ through a single boundary-norm estimate for \(D^3H\), where \(D=t\frac{d}{dt}\). Rigorous interval arithmetic verifies the resulting finite collection of inequalities, yielding the claimed upper bound on \(K_G\).



% The computational verification is reduced to a simple criterion on the coefficients of the correlation function. 
% \rnote{What computational verification; this is a big jump from above. This is almost assuming familiarity with our previous paper.}
% Write
% \(H(t)=b_1t+\sum_{m\geq 3}b_mt^m\), and let
% \(\Delta_H:=\sum_{m\geq 3}|b_m|\) denote its total nonlinear coefficient mass. We show that
% \[
%     \gamma+\Delta_H<b_1
%     \qquad\Longrightarrow\qquad
%     M_H(\gamma)<1.
% \]
% Thus, rather than computing the inverse series \(H^{-1}\), it is enough to certify a lower bound on the linear coefficient \(b_1\) and an upper bound on \(\Delta_H\). We reduce these two estimates to a finite computer-assisted calculation. The initial coefficients of \(H\) are enclosed directly using one-dimensional Gaussian integrals, while the remaining coefficient tail is bounded through a single boundary-norm estimate for \(D^3H\), where \(D=t\frac{d}{dt}\). Rigorous interval arithmetic verifies the resulting finite collection of inequalities, yielding the claimed upper bound on \(K_G\).


\subsection{The lower bound}

\paragraph{Mixed Krivine schemes are optimal.}
Naor and Regev~\cite{naor2014krivine} proved that mixed Krivine schemes are asymptotically optimal. We use a quantitative consequence of their construction: there exist limiting mixed correlation functions \(H_k\) with admissible inverse-majorant parameters \(\gamma_k\) satisfying
\[
    \gamma_k\longrightarrow\frac{\pi}{2K_G}.
\]
Consequently, any universal upper bound on the admissible parameter of a Krivine correlation function gives a lower bound on \(K_G\).


\paragraph{The affine coefficient inequality.}
To obtain this upper bound, we ask which Taylor coefficients can actually arise from a Krivine correlation function. For odd sign functions \(f,g:\mathbb R^d\to\{-1,1\}\), write
\[
    H_{f,g}(t)=b_1t+b_3t^3+\cdots.
\]

Surprisingly, we show that there is the clean constraint on the first two coefficients:
\begin{equation}
    b_3\geq 2b_1-\frac{11}{6}.
    \label{eq:affine-coefficient-inequality}
\end{equation}
Its affine form is important, since it is preserved under mixtures and coefficientwise limits. 

Pushing the affine inequality to the inverse series $H^{-1}$ yields bounds on the inverse coefficients $a_1$ and $a_3$, which in turn gives $\gamma \leq \frac{11}{12}$. This results in the desired $K_G \geq \frac{6\pi}{11}$ .


% It also gives the desired barrier on the inverse majorant. Indeed, if
% \[
%     H^{-1}(z)=a_1z+a_3z^3+\cdots,
% \]
% then
% \[
%     a_1=\frac{1}{b_1},
%     \qquad
%     a_3=-\frac{b_3}{b_1^4}.
% \]
% Combining these identities with \eqref{eq:affine-coefficient-inequality}, the bound \(b_1\leq 1\), and the first two terms of the inverse majorant shows that every admissible parameter satisfies
% \[
%     \gamma\leq\frac{11}{12}.
% \]
% Since the asymptotically optimal family above has
% \(\gamma_k\to\pi/(2K_G)\), this yields
% \[
%     K_G\geq\frac{6\pi}{11}.
% \]

% \paragraph{Agreement and disagreement.}
% We prove \eqref{eq:affine-coefficient-inequality} using the Hermite expansions of \(f\) and \(g\). Let \(P_iF\) denote the projection of a function \(F\) onto Hermite degree \(i\). Mehler's identity gives
% \[
%     b_1=\frac{\pi}{2}\langle P_1f,P_1g\rangle,
%     \qquad
%     b_3=\frac{\pi}{2}\langle P_3f,P_3g\rangle.
% \]
% Writing
% \[
%     \nu:=\mathbb E|Z|=\sqrt{\frac{2}{\pi}},
% \]
% the affine inequality is therefore equivalent to
% \[
%     2\langle P_1f,P_1g\rangle
%     -\langle P_3f,P_3g\rangle
%     \leq \frac{11}{6}\nu^2.
% \]

% The useful change of variables is
% \[
%     h:=\frac{f+g}{2},
%     \qquad
%     k:=\frac{f-g}{2}.
% \]
% Here \(h\) is supported where \(f\) and \(g\) agree, while \(k\) is supported where they disagree. Moreover,
% \[
%     h,k\in\{-1,0,1\},
%     \qquad
%     h^2+k^2=1.
% \]
% Substituting \(f=h+k\) and \(g=h-k\), the left-hand side above becomes
% \[
%     2\|P_1h\|^2-2\|P_1k\|^2
%     -\|P_3h\|^2+\|P_3k\|^2.
% \]
% After dropping the nonpositive term \(-2\|P_1k\|^2\), it is enough to prove
% \begin{equation}
%     2\|P_1h\|^2-\|P_3h\|^2+\|P_3k\|^2
%     \leq \frac{11}{6}\nu^2.
%     \label{eq:agreement-disagreement-reduction}
% \end{equation}
% The first two terms involve the agreement function \(h\), while the final term measures the degree-three contribution of the disagreement function \(k\). We bound these two contributions separately.









% \paragraph{The affine coefficient inequality.}
% We can ask 



% Pushing the affine inequality to the inverse series $H^{-1}$ yields bounds on the inverse coefficients $a_1$ and $a_3$, which in turn gives $\gamma \leq \frac{11}{12}$. This results in the desired $K_G \leq \frac{6\pi}{11}$ .




% \paragraph{The affine coefficient inequality.}
% For a square-integrable function \(F\) on Gaussian space, let \(P_iF\) denote its orthogonal projection onto the Hermite polynomials of total degree \(i\). If
% \(H_{f,g}(t)=b_1t+b_3t^3+\cdots\) is the normalized correlation function of \(f\) and \(g\), then the basic orthogonality properties of Hermite polynomials give
% \(b_1=\frac{\pi}{2}\langle P_1f,P_1g\rangle\) and
% \(b_3=\frac{\pi}{2}\langle P_3f,P_3g\rangle\).
% Our main coefficient estimate is
% \begin{equation}
%     b_3\geq 2b_1-\frac{11}{6}.
%     \label{eq:affine-coefficient-inequality}
% \end{equation}
% Since \(\nu:=\mathbb E|Z|=\sqrt{2/\pi}\), this is equivalent to
% \[
%     2\langle P_1f,P_1g\rangle-\langle P_3f,P_3g\rangle
%     \leq \frac{11}{6}\nu^2.
% \]

% \rnote{Talking about $P_i F$ here seems a bit odd. I don't think that'll help in understanding yet. It might help to motivate the general approach. Something like ... We can ask what sets of coeffecients $(b_1, b_3, \cdots)$ can in fact be realized as Taylor-coefficients of a correlation function. Surprisingly, we show that there are some non-trivial clean constraints: ...}

\paragraph{Hermite expansion.} The approach to proving this is by looking at Hermite expansions of $f, g$. For a square-integrable function \(F\) on Gaussian space, let \(P_iF\) denote its orthogonal projection onto the Hermite polynomials of total degree \(i\).  If
\(H_{f,g}(t)=b_1t+b_3t^3+\cdots\) is the normalized correlation function of \(f\) and \(g\), then the basic orthogonality properties of Hermite polynomials give
\(b_1=\frac{\pi}{2}\langle P_1f,P_1g\rangle\) and
\(b_3=\frac{\pi}{2}\langle P_3f,P_3g\rangle\).
Our affine coefficient inequality is
\begin{equation*}
    b_3\geq 2b_1-\frac{11}{6}.
\end{equation*}
Since \(\nu:=\mathbb E|Z|=\sqrt{2/\pi}\), this is equivalent to
\[
    2\langle P_1f,P_1g\rangle-\langle P_3f,P_3g\rangle
    \leq \frac{11}{6}\nu^2.
\]

\paragraph{Agreement and disagreement.}
In order to prove the inequality, we make a key substitution. The substitution is to write
\(h:=(f+g)/2\) and \(k:=(f-g)/2\). The function \(h\) can be thought of as the \textit{agreement part}, where \(f\) and \(g\) agree, while \(k\) is the corresponding \textit{disagreement part}. Moreover,
\(h,k\in\{-1,0,1\}\) and \(h^2+k^2=1\).

% \rnote{This section seems like its jumping. After stating the inequality clearly for $b_3, b_1$, you can then say we approach proving this by looking at Hermite expansions of $f,g$. In particualr $b_3 = blah $ .. }
Substituting \(f=h+k\) and \(g=h-k\) into the preceding inequality, and discarding a nonpositive term, it suffices to prove
\begin{equation}
    2\|P_1h\|^2-\|P_3h\|^2+\|P_3k\|^2
    \leq \frac{11}{6}\nu^2.
    \label{eq:agreement-disagreement-reduction}
\end{equation}
The first two terms depend only on the agreement function \(h\), while the final term measures the degree-three contribution of the disagreement function \(k\). We bound these two parts separately.


\paragraph{Agreement terms.}
Rotate coordinates so that \(P_1h\) lies in a single Gaussian direction \(S\), and write
\(x:=\|P_1h\|/\nu\in[0,1]\). The positive agreement term is then
\(2\|P_1h\|^2=2\nu^2x^2\).

To control the helpful negative term \(-\|P_3h\|^2\), average \(h\) over all coordinates transverse to \(S\). Namely, writing the Gaussian input as \((S,T)\), set
\(A(s):=\mathbb E_T[h(s,T)]\). Then \(|A|\leq1\) and
\(\mathbb E[A(S)S]=\nu x\). A one-dimensional rearrangement inequality shows that this prescribed linear moment forces a cubic Hermite moment:
\[
    \|P_3h\|^2
    \geq
    \frac{\nu^2}{6}\Delta_+(x)^2
\]
for an explicit function \(\Delta\). Thus, a large linear agreement term necessarily creates a compensating cubic agreement term.

\paragraph{Disagreement term.}
For each fixed transverse vector \(t\), consider the one-dimensional slice
\(u_t(s):=k(s,t)\). Since \(k\in\{-1,0,1\}\), each \(u_t\) is a ternary function on the Gaussian line. The degree-three Hermite decomposition implies
\[
    \|P_3k\|^2
    \leq
    \mathbb E_T V(u_T),
    \qquad
    V(u):=
    \sum_{j=0}^3
    \bigl(\mathbb E[u(Z)\psi_j(Z)]\bigr)^2.
\]
We control these first four Hermite moments through the weighted support
\(\beta(u):=\mathbb E[|Z|u(Z)^2]\). The key one-dimensional estimate is
\begin{equation}
    V(u)\leq 3\nu\beta(u)-\beta(u)^2.
    \label{eq:one-dimensional-support-inequality}
\end{equation}
Conceptually, this says that a ternary function cannot have both a support concentrated cheaply near the origin and a large projection onto all Hermite degrees at most three. The inequality admits an exact finite-dimensional dual reduction. After introducing a support price \(c\), the optimizing ternary function has the threshold form
\(
    u(z)=\operatorname{sgn}(p(z))
    \mathbf 1_{\{|p(z)|\geq c|z|\}},
\)
where \(p\) is a unit cubic Hermite polynomial. The resulting one-dimensional Gaussian integral inequalities are verified using outward-rounded interval arithmetic. Finally, \(h^2+k^2=1\) relates the average weighted support of the slices \(u_T\) to the agreement parameter \(x\). Combining this relation with
\eqref{eq:one-dimensional-support-inequality} gives an explicit upper bound on
\(\|P_3k\|^2\) in terms of \(x\). Together with the agreement estimate, this reduces
\eqref{eq:agreement-disagreement-reduction} to an elementary scalar inequality on
\(x\in[0,1]\), proving \eqref{eq:affine-coefficient-inequality}.




% \section{Technical Overview}

% Here, we provide an exposition of the main ideas in our work. 

% \paragraph{Classical Krivine schemes.} ... 

% \paragraph{Limiting Krivine schemes.} ... 

% \paragraph{Approximation of limiting schemes with classical schemes.} ...

% \paragraph{The cubic-quintic scheme.} ...

% \paragraph{Deriving the upper bound.} ...

% \paragraph{Krivine schemes are optimal.} ... 

% \paragraph{Affine inequality.} ... 

% \paragraph{Deriving the lower bound.} ... 

% \paragraph{Technical overview of the Lower Bound.} We prove the following theorem:

% \begin{theorem}\label{thm:lower-bound}
% The Grothendieck constant satisfies
% \[
%         K_G\ge\frac{6\pi}{11} \approx 1.7136.
% \]
% \end{theorem}

% A sketch of the proof is as follows. 

% First, recall the inverse-series step in Krivine's rounding argument. If
% \[
%         H^{-1}(z)=\sum_{j\ge0}a_{2j+1}z^{2j+1},
% \]
% define the absolute-coefficient majorant $M_H(s):=\sum_{j\ge0}|a_{2j+1}|s^{2j+1}$. The standard tensor preprocessing gives the implication
% \begin{equation}
%         M_H(\gamma)\le1
%         \quad\Longrightarrow\quad
%         K_G\le\frac{\pi}{2\gamma}.
%         \label{eq:inverse-majorant-bound}
% \end{equation}
% We refer to~\eqref{eq:inverse-majorant-bound} as the \emph{inverse-majorant bound}.

% The conceptual starting point is the theorem of Naor and Regev that Krivine schemes are optimal \cite{naor2014krivine}. By reworking the construction in their proof, we show that the inverse-majorant bound is itself asymptotically optimal after allowing mixtures and increasing dimension: there are limiting mixed correlation functions with admissible parameters converging to $\pi/(2K_G)$. This transfer is established in Section~\ref{sec:majorant_optimal}.

% It therefore remains to prove a universal upper bound on such admissible parameters. For arbitrary odd measurable sign functions $f,g:\R^d\to\{\pm1\}$, write
% \[
%         H_{f,g}(t)=b_1t+b_3t^3+b_5t^5+\cdots.
% \]
% The technical core of the argument is the affine inequality
% \[
%         b_3\ge 2b_1-\frac{11}{6}.
% \]
% Section~\ref{sec:preliminaries} sets up the Gaussian--Hermite formulation, and Section~\ref{sec:affine-strip} proves this inequality by decomposing $f$ and $g$ into their agreement and disagreement parts and reducing the required estimates to one-dimensional inequalities. Its linearity in $(b_1,b_3)$ is essential, because it makes the inequality invariant under mixtures and coefficientwise limits.

% Section~\ref{sec:strip_majorant} then shows that the affine inequality controls the inverse-majorant bound: every admissible value must satisfy $\gamma\le11/12$. Applying this barrier to the asymptotically optimal sequence constructed in Section~\ref{sec:majorant_optimal} gives
% \[
%         \frac{\pi}{2K_G}\le\frac{11}{12},
% \]
% which is exactly Theorem~\ref{thm:lower-bound}.


% \paragraph{Technical overview of the Upper Bound}. (Give references to sections of the technical part)

% - Recall classical Krivine schemes... (standard correlation)

% - Define Limiting Krivine schemes... (arbitrary correlation)

% - Show that these schemes can be arbitrarily approximated by classical (finite) Krivine schemes, in a specific sense

% - Therefore inverse-majorant bound on Limiting Krivine schemes translates to upper bounds on $K_G$, so this is a new search space

% - We then study a specific instantiation of limiting Krivine schemes: the cubic--quintic scheme...

% - We show how to derive the inverse-majorant bound on this by bounding the tail of the taylor expansion of $H$... 


